% =====================================================================
%  CLASE 02 -- Hidráulica de pozos en régimen permanente
% =====================================================================
\clase{02}{Aguas subterráneas: hidráulica de pozos en régimen permanente}

\section{Recordando\ldots}
\begin{itemize}
  \item \textbf{Acuífero libre:} cercano a la superficie terrestre; el agua
        subterránea tiene superficie abierta a la atmósfera.
  \item \textbf{Acuífero artesiano o confinado:} confinado entre estratos
        impermeables, por ende, bajo presión.
  \item \textbf{Ley de Darcy (1856)}, válida para flujo laminar ($Re<4$):
  \[
  Q=-K\cdot A\cdot\left(\frac{\dd h}{\dd l}\right)=-K\cdot A\cdot i
  \]
\end{itemize}

\section{Propiedades de los acuíferos}

\subsection{Coeficiente de almacenamiento específico $S_s$ [1/L]}
Es la cantidad de agua, por unidad de volumen, que es almacenada o liberada
debido a la compresibilidad del esqueleto mineral y del agua en los poros,
debido a un cambio unitario en el nivel de agua en el acuífero [1/L].
\begin{formula}
$S_s=g\rho(\alpha+\eta\beta)$
\end{formula}
donde $g$: aceleración de gravedad [m/s$^2$]; $\rho$: densidad del agua
[kg/m$^3$]; $\eta$: porosidad del medio; $\alpha$: compresibilidad de la
matriz sólida del acuífero [m$^2$/N]; $\beta$: compresibilidad del agua [m$^2$/N].

\subsection{Coeficiente de almacenamiento $S$}
Volumen de agua, por unidad de área y cambio de altura, que la unidad
permeable absorbe o libera.

\paragraph{Acuífero confinado de espesor $m$:}
\begin{formula}
$S=S_s\,m\qquad\qquad V_{\text{drenado}}=-S\cdot A\cdot\Delta h$\\[2pt]
Acuífero confinado: $10^{-4}<S<0.005$
\end{formula}

\begin{figure}[H]
\centering
\begin{tikzpicture}[font=\small]
  \fill[gray!30] (0,0) rectangle (7,0.8);  \node at (1,0.4) {Acuitardo};
  \fill[cyan!40] (0,0.8) rectangle (7,2.4); \node at (1,1.6) {Acuífero};
  \fill[gray!30] (0,2.4) rectangle (7,3.3); \node at (1,2.85) {Acuitardo};
  \draw[thick] (0,3.6) -- (7,3.4);
  \draw[thick,dashed] (0,3.3) -- (7,3.1);
  \draw[decorate,decoration={brace,amplitude=3pt}] (7.1,3.4) -- (7.1,3.1)
    node[midway,right,align=left,font=\scriptsize]{Descenso unitario de la\\superficie piezométrica};
  % prisma unitario
  \draw[dashed] (4,0.8) rectangle (4.8,2.4);
  \draw[dashed] (4,2.4) -- (4.3,2.65) -- (5.1,2.65) -- (4.8,2.4);
  \draw[dashed] (5.1,2.65) -- (5.1,1.05) -- (4.8,0.8);
  \node[font=\scriptsize] at (4.4,2.85) {1};\node[font=\scriptsize] at (4.95,2.85) {1};
  \draw[cota] (5.6,0.8) -- (5.6,2.4) node[midway,right]{$m$};
\end{tikzpicture}
\caption{Coeficiente de almacenamiento en un acuífero confinado.}
\end{figure}

\paragraph{Acuífero libre:}
\begin{itemize}
  \item En acuíferos libres también se conoce como \textbf{rendimiento
        específico} ($S=S_y$).
  \item Coincide con la porosidad efectiva ($S_y=n_e$).
\end{itemize}
\begin{formula}
$V_{\text{drenado}}=-S\cdot A\cdot\Delta h$\qquad Acuífero libre: $0.05<S<0.2$
\end{formula}

\begin{table}[H]
\centering
\begin{tabular}{lc}
\toprule
\textbf{Material} & $S_y$\\
\midrule
Arcilla & 0.00 -- 0.05\\
Limo & 0.03 -- 0.19\\
Arena fina & 0.10 -- 0.32\\
Arena media & 0.15 -- 0.32\\
Arena gruesa & 0.20 -- 0.35\\
Grava & 0.14 -- 0.35\\
\bottomrule
\end{tabular}
\caption{Valores típicos de rendimiento específico (Johnson, 1967).}
\end{table}

\begin{figure}[H]
\centering
\begin{tikzpicture}[font=\scriptsize]
  \foreach \X/\tit/\lab in {0/{Acuífero confinado}/S, 5.5/{Acuífero libre}/m_e}{
  \begin{scope}[xshift=\X cm]
    \draw (0,0) rectangle (1.6,5);
    \draw (1.6,0) -- (2.1,0.4) -- (2.1,5.4) -- (0.5,5.4) -- (0,5);
    \draw (1.6,5) -- (2.1,5.4);
    \node[rotate=15] at (1.05,5.2) {1 m$^2$};
    \fill[pattern=dots,pattern color=gray] (0,0) rectangle (1.6,2.6);
    \draw[dashed] (0,3.6) -- (1.6,3.6) (0,2.6) -- (1.6,2.6);
    \draw[cota] (1.8,2.6) -- (1.8,3.6) node[midway,right]{1 m};
    \draw (-0.8,4.8) rectangle (-0.3,5.4); \node at (-0.55,5.1) {$\lab$};
    \node[align=center] at (0.8,-0.5) {\tit};
  \end{scope}}
  \node[align=left] at (-0.3,6.2) {Extrayendo un volumen $S$ desciende\\la superficie piezométrica 1 metro};
  \node[align=left] at (5.4,6.2) {Extrayendo un volumen $m_e$ desciende\\la superficie freática 1 metro};
\end{tikzpicture}
\caption{Interpretación física del coeficiente de almacenamiento en acuíferos
confinado y libre.}
\end{figure}

\section{Ecuación de conservación de masa en un medio saturado}

Para un acuífero confinado, el balance de masa sobre un elemento diferencial,
asumiendo fluido incompresible y flujo sólo en la dirección $x$, queda:
\[
\rho\left(v_x-\left(v_x+\frac{\partial v_x}{\partial x}\,\dd x\right)\right)\dd y\,\dd z
=\frac{\partial M}{\partial t}
\quad\Longrightarrow\quad
\boxed{-\rho\,\frac{\partial v_x}{\partial x}\,\dd x\,\dd y\,\dd z=\frac{\partial M}{\partial t}}
\]

\begin{figure}[H]
\centering
\begin{tikzpicture}[font=\small]
  \fill[gray!25] (0,0) rectangle (2.5,2.5);
  \fill[gray!45] (2.5,0) -- (3.3,0.6) -- (3.3,3.1) -- (2.5,2.5) -- cycle;
  \fill[gray!15] (0,2.5) -- (0.8,3.1) -- (3.3,3.1) -- (2.5,2.5) -- cycle;
  \draw[flecha,cyan] (-1.2,1.25) -- (-0.1,1.25) node[pos=0,above right]{$v_x$};
  \draw[flecha,cyan] (2.9,1.6) -- (4,1.6) node[right,black]{$v_x+\dfrac{\partial v_x}{\partial x}\dd x$};
  \node[below] at (0,0) {$x$}; \node[below] at (2.5,0) {$x+\Delta x$};
  \node[below] at (1.25,0) {$\Delta x$};
  \node[right] at (3.3,0.3) {$\Delta y$}; \node[right] at (3.3,2.9) {$\Delta z$};
\end{tikzpicture}
\caption{Elemento diferencial para el balance de masa.}
\end{figure}

Considerando que
\[
M=\rho V_{\text{drenado}}=-\rho S_s\,\dd z\,A\,\Delta h,\qquad A=\dd x\,\dd y
\quad\Longrightarrow\quad
\frac{\partial M}{\partial t}=-\rho S_s\frac{\partial h}{\partial t}\,\dd x\,\dd y\,\dd z
\]
Si se asume flujo en las tres direcciones:
\begin{formula}
$\displaystyle -\left(\frac{\partial v_x}{\partial x}+\frac{\partial v_y}{\partial y}+\frac{\partial v_z}{\partial z}\right)=-S_s\frac{\partial h}{\partial t}$
\end{formula}
\emph{¿Y de dónde obtenemos las velocidades?} Recordando que Darcy (1856)
estableció que las velocidades de escurrimiento en medios porosos están
dadas por:
\[
v_x=-K_x\frac{\partial h}{\partial x},\qquad
v_y=-K_y\frac{\partial h}{\partial y},\qquad
v_z=-K_z\frac{\partial h}{\partial z}
\]
De este modo se deduce la \textbf{ecuación general de escurrimiento en un
medio saturado confinado}:
\begin{formula}
$\displaystyle
\frac{\partial}{\partial x}\left(K_x\frac{\partial h}{\partial x}\right)+
\frac{\partial}{\partial y}\left(K_y\frac{\partial h}{\partial y}\right)+
\frac{\partial}{\partial z}\left(K_z\frac{\partial h}{\partial z}\right)=S_s\frac{\partial h}{\partial t}$\\[4pt]
$S_s$: almacenamiento específico, $S/m$
\end{formula}

\section{Ecuaciones generales}
\subsection{Casos particulares}
\begin{itemize}
  \item Flujo transiente en medio homogéneo e isotrópico:
  \[\frac{\partial^2h}{\partial x^2}+\frac{\partial^2h}{\partial y^2}+\frac{\partial^2h}{\partial z^2}=\frac{S_s}{K}\frac{\partial h}{\partial t}\]
  \item Flujo permanente, anisotrópico, homogéneo:
  \[\frac{\partial}{\partial x}\left(K_x\frac{\partial h}{\partial x}\right)+\frac{\partial}{\partial y}\left(K_y\frac{\partial h}{\partial y}\right)+\frac{\partial}{\partial z}\left(K_z\frac{\partial h}{\partial z}\right)=0\]
  \item Flujo permanente, isotrópico, homogéneo (\textbf{ecuación de Laplace}):
\end{itemize}
\begin{formula}
$\displaystyle\frac{\partial^2h}{\partial x^2}+\frac{\partial^2h}{\partial y^2}+\frac{\partial^2h}{\partial z^2}=0$\qquad\textbf{Laplace}
\end{formula}

\subsection{Ecuación de Laplace en coordenadas cilíndricas}
Flujo permanente, isotrópico y homogéneo:
\begin{formula}
$\displaystyle\frac{1}{r}\frac{\partial h}{\partial r}+\frac{\partial^2h}{\partial r^2}+\frac{1}{r^2}\frac{\partial^2h}{\partial\theta^2}+\frac{\partial^2h}{\partial z^2}=0$
\end{formula}

\subsection{Métodos de solución}
\begin{enumerate}
  \item Solución directa exacta.
  \item Integración analítica simplificada.
  \item Redes de flujo.
  \item Modelos analógicos.
  \item Modelos numéricos.
\end{enumerate}

\paragraph{¿Para qué hacemos esto?} ¿Por qué es importante saber la dinámica
de la altura del nivel freático en el tiempo? Porque al bombear un pozo se
genera un \emph{cono de depresión} que determina cuánto desciende la napa en
los alrededores, y con ello el caudal que es posible extraer.

\section{Hidráulica de pozos}
\subsection{Supuestos básicos}
\begin{itemize}
  \item Los acuíferos se considerarán de propiedades homogéneas e isotrópicas,
        de extensión infinita ($S$, $m$ constantes).
  \item El tipo de acuífero se supone invariante en el espacio y el tiempo,
        i.e., es libre, confinado o semiconfinado.
  \item Se supone que la superficie piezométrica es horizontal antes del
        inicio de operación del pozo.
  \item El flujo será radial concéntrico hacia el pozo, y tendrá un radio de
        influencia por sobre el cual se supondrá despreciable el descenso de la napa.
  \item No existen pérdidas de energía por roce al escurrir agua hacia el pozo.
  \item No se considera (por ahora) superposición de pozos.
\end{itemize}

% ---------------------------------------------------------------------
%  Macro para dibujar un pozo en acuífero (confinado=1 / libre=0)
% ---------------------------------------------------------------------
\newcommand{\dibujopozo}[1]{%
\begin{tikzpicture}[font=\small,scale=0.95]
  % base impermeable
  \draw[thick] (-4,0) -- (4,0);
  \foreach \x in {-3.9,-3.6,...,3.9} \draw (\x,0) -- (\x-0.2,-0.25);
  % pozo
  \draw[thick] (-0.35,0.1) -- (-0.35,4.6) (0.35,0.1) -- (0.35,4.6);
  \draw[thick] (-4,4.6) -- (-0.35,4.6) (0.35,4.6) -- (4,4.6);
  \fill[azulusm] (-0.1,4.5) rectangle (0.1,5.1);
  \fill[azulusm] (-0.25,5.1) -- (0.25,5.1) -- (0,5.45) -- cycle;
  \node at (0.5,5.35) {$Q$};
  % superficie original y cono
  \draw[dashed,cyan!70!blue,thick] (-4,3.9) -- (4,3.9);
  \node[azulusm,font=\scriptsize\bfseries,anchor=south west] at (-4.2,3.95) {Superficie piezométrica original};
  \draw[cyan!70!blue,thick] (-4,3.85) .. controls (-2,3.7) and (-1,3.2) .. (-0.35,2.4) -- (0.35,2.4)
      .. controls (1,3.2) and (2,3.7) .. (4,3.85);
  \pic at (0,2.42) {nivelagua};
  \ifnum#1=1
    % techo impermeable del acuífero confinado
    \draw[thick] (-4,2.0) -- (-0.35,2.0) (0.35,2.0) -- (4,2.0);
    \foreach \x in {-3.9,-3.6,...,-0.5,0.6,0.9,...,3.9} \draw (\x,2.0) -- (\x+0.15,2.25);
    \draw[cota] (-3.5,0) -- (-3.5,2.0) node[midway,left]{$m$};
    \node[azulusm,font=\scriptsize\bfseries,align=center] at (-2.6,3.1) {Superficie\\piezométrica};
    \def\ytop{1.9}
  \else
    \node[azulusm,font=\scriptsize,align=center] at (-2.9,2.9) {$\dfrac{\partial h}{\partial s}\approx\dfrac{\partial h}{\partial r}$\\Dupuit--Forchheimer};
    \def\ytop{2.2}
  \fi
  % flechas de flujo
  \foreach \y in {0.3,0.6,...,\ytop}{
    \draw[->,azulusm] (-1.2,\y) -- (-0.5,\y);
    \draw[->,azulusm] (1.2,\y) -- (0.5,\y);}
  % cotas
  \draw[cota] (-0.15,0) -- (-0.15,2.4) node[pos=0.8,left]{$h_0$};
  \draw[->] (0,1.2) -- (0.35,1.2) node[right,font=\scriptsize]{$r_0$};
  \draw[cota] (2.2,0) -- (2.2,3.62) node[midway,right]{$h$};
  \draw[cota] (3.8,0) -- (3.8,3.9) node[midway,left]{$H$};
  \draw[->] (0,-0.5) -- (2.2,-0.5) node[midway,below]{$r$};
  \draw[->] (0,-1.1) -- (3.9,-1.1) node[midway,below,align=center]{$R$\\[-2pt]\scriptsize\color{azulusm}\bfseries Radio de influencia};
\end{tikzpicture}}

\subsection{Pozo de penetración total en acuífero confinado}
\begin{itemize}
  \item Extensión lateral infinita.
  \item Pozo bombea caudal constante.
  \item Líneas de flujo radiales y convergentes hacia el pozo.
  \item Equipotenciales son círculos concéntricos.
\end{itemize}

\begin{figure}[H]
\centering
\dibujopozo{1}
\caption{Pozo de penetración total en acuífero confinado.}
\end{figure}

\begin{nota}[title={Deducción (ecuación de Thiem)}]
Por continuidad, a través de un cilindro de radio $r$ y altura $m$:
\[
Q=2\pi r\,m\,K\frac{\dd h}{\dd r}\;\Rightarrow\;\int_h^H\dd h=\frac{Q}{2\pi Km}\int_r^R\frac{\dd r}{r}
\]
\end{nota}
\begin{formula}
$\displaystyle\Delta h=H-h=\frac{Q}{2\pi T}\ln\!\left(\frac{R}{r}\right)$\\[4pt]
$R$: radio de influencia;\quad $T=Km$: transmisibilidad
\end{formula}

\subsection{Pozo de penetración total en acuífero libre}
\begin{itemize}
  \item Pozo bombea caudal constante.
  \item Líneas de flujo radiales y convergentes hacia el pozo.
  \item Equipotenciales son círculos concéntricos.
  \item Flujo horizontal.
  \item Gradiente hidráulico no varía con la profundidad.
  \item Se cumple ley hidrostática.
\end{itemize}
Estas son las hipótesis de \textbf{Dupuit--Forchheimer}:
$\dfrac{\partial h}{\partial s}\approx\dfrac{\partial h}{\partial r}$.

\begin{figure}[H]
\centering
\dibujopozo{0}
\caption{Pozo de penetración total en acuífero libre.}
\end{figure}

\begin{nota}[title={Deducción (Dupuit)}]
Ahora el área de flujo es $2\pi r h$:
\[
Q=2\pi r\,h\,K\frac{\dd h}{\dd r}\;\Rightarrow\;\int_h^H2h\,\dd h=\frac{Q}{\pi K}\int_r^R\frac{\dd r}{r}
\]
\end{nota}
\begin{formula}
$\displaystyle H^2-h^2=\frac{Q}{\pi K}\ln\!\left(\frac{R}{r}\right)$\qquad $R$: radio de influencia
\end{formula}

\paragraph{Aproximaciones.} Si el espesor inicial de la napa saturada es muy
grande en comparación con las depresiones cerca del pozo, la ecuación se
simplifica y es similar a la de un acuífero confinado:
\begin{itemize}
  \item Si $\Delta h<0.1\,H$:
  \[\Delta h=\frac{Q}{2\pi KH}\ln\frac{R}{r}\]
  \item Si $\Delta h<0.5\,H$ (\textbf{corrección de Jacob}):
  \[\Delta h'=\Delta h-\frac{(\Delta h)^2}{2H}=\frac{Q}{2\pi KH}\ln\frac{R}{r}\]
\end{itemize}

\section{Acuíferos libres de gran potencia}

\begin{figure}[H]
\centering
\begin{subfigure}{0.48\textwidth}\centering
\begin{tikzpicture}[scale=0.8,font=\small]
  \draw[thick] (-2.5,4) -- (-0.4,4) (0.4,4) -- (2.5,4);
  \foreach \x in {-2.4,-2.0,...,-0.6,0.6,1.0,...,2.4} \draw (\x,4) -- (\x+0.25,4.25);
  \draw[thick] (-0.4,4.6) -- (-0.4,1.3) (0.4,4.6) -- (0.4,1.3);
  \draw[cyan,thick] (-2.5,3.2) -- (2.5,3.2); \pic at (-2.1,3.22) {nivelagua};
  \shade[ball color=azulusm!70] (0,1) circle (0.8);
  \foreach \a in {0,45,...,315} \draw[->,azulusm,thick] (\a:2.2) ++(0,1) -- ++(\a+180:1.1);
  \node at (1.6,4.5) {$H>R$};
\end{tikzpicture}
\caption{Pozo en acuífero libre de gran potencia\\(escurrimiento esférico).}
\end{subfigure}\hfill
\begin{subfigure}{0.48\textwidth}\centering
\begin{tikzpicture}[scale=0.8,font=\small]
  \draw[thick] (-2.5,4) -- (-0.4,4) (0.4,4) -- (2.5,4);
  \foreach \x in {-2.4,-2.0,...,-0.6,0.6,1.0,...,2.4} \draw (\x,4) -- (\x+0.25,4.25);
  \draw[thick] (-0.4,4.6) -- (-0.4,2.6) (0.4,4.6) -- (0.4,2.6);
  \draw[cyan,thick] (-2.5,2.8) -- (2.5,2.8); \pic at (-2.1,2.82) {nivelagua};
  \fill[azulusm!70] (-0.7,2.8) arc (180:360:0.7) -- cycle;
  \foreach \a in {0,-45,-90,-135,-180} \draw[->,azulusm,thick] (\a:2.0) ++(0,2.8) -- ++(\a+180:1.1);
\end{tikzpicture}
\caption{Pozo somero en acuífero libre de gran potencia (escurrimiento semi-esférico).}
\end{subfigure}
\end{figure}

\begin{formula}
Escurrimiento esférico: $\displaystyle\Delta h=H-h=\frac{Q}{4\pi K}\left(\frac1r-\frac1R\right)$\\[6pt]
Escurrimiento semi-esférico: $\displaystyle\Delta h=H-h=\frac{Q}{2\pi K}\left(\frac1r-\frac1R\right)$,
\quad $r_0\approx0.6\sim8$ [m]
\end{formula}

\section{Pozo de penetración parcial en acuífero confinado}
Sea $b$ la longitud de penetración del pozo en un acuífero confinado de
espesor $m$.

\begin{figure}[H]
\centering
\begin{tikzpicture}[font=\small,scale=0.9]
  \draw[thick] (-4,0) -- (4,0);
  \foreach \x in {-3.9,-3.6,...,3.9} \draw (\x,0) -- (\x-0.2,-0.25);
  \draw[thick] (-4,3) -- (-0.35,3) (0.35,3) -- (4,3);
  \foreach \x in {-3.9,-3.6,...,-0.5,0.6,0.9,...,3.9} \draw (\x,3) -- (\x+0.15,3.25);
  \draw[thick] (-0.35,1.2) -- (-0.35,4.6) (0.35,1.2) -- (0.35,4.6);
  \draw[thick] (-4,4.6) -- (-0.35,4.6) (0.35,4.6) -- (4,4.6);
  \fill[azulusm] (-0.1,4.5) rectangle (0.1,5.1);
  \fill[azulusm] (-0.25,5.1) -- (0.25,5.1) -- (0,5.45) -- cycle; \node at (0.5,5.35) {$Q$};
  \draw[dashed,cyan!70!blue,thick] (-4,4.2) .. controls (-2,4.0) and (-1,3.8) .. (-0.35,3.6) -- (0.35,3.6) .. controls (1,3.8) and (2,4.0) .. (4,4.2);
  \foreach \y in {1.5,1.8,...,2.8}{\draw[->,azulusm] (-1.2,\y) -- (-0.5,\y);\draw[->,azulusm] (1.2,\y) -- (0.5,\y);}
  \foreach \a in {-60,-90,-120}\draw[->,azulusm] (\a:0.9) ++(0,1.2) -- ++(\a+180:0.5);
  \draw[cota] (-3.5,0) -- (-3.5,3) node[midway,left]{$m$};
  \draw[cota] (1.6,1.2) -- (1.6,3) node[midway,right]{$b$};
  \draw[->] (0,-0.5) -- (2.5,-0.5) node[midway,below]{$r$};
\end{tikzpicture}
\caption{Pozo de penetración parcial en acuífero confinado.}
\end{figure}

\paragraph{Solución de Nasberg} (válida para $b/m\geq0.3$):
\[
h_1-h_2=\frac{Q}{2\pi Km}\left(\ln\!\left(\frac{r_1}{r_2}\right)+\operatorname{senh}^{-1}\!\left(\frac{m}{r_1}\right)-\operatorname{senh}^{-1}\!\left(\frac{m}{r_2}\right)-\frac{m}{b}\left(\operatorname{senh}^{-1}\!\left(\frac{b}{r_1}\right)-\operatorname{senh}^{-1}\!\left(\frac{b}{r_2}\right)\right)\right)
\]

\paragraph{Solución de Girinsky} (válida para $b/m<0.3$):
\[
h_1-h_2=\frac{Q}{2\pi K}\left(\frac1b\ln\frac{1.6\,b}{r_2}-\operatorname{arcsenh}\frac{b}{r_1}\right)
\]

\section{Estimación del radio de influencia $R$}
\paragraph{En terreno}
\begin{itemize}
  \item Mediciones de depresiones en función de la distancia al pozo.
  \item Determinar $R$ cuando la depresión se hace cero (gráfico $\Delta h$ vs.\ $\ln r$).
\end{itemize}
\begin{figure}[H]
\centering
\begin{tikzpicture}[font=\small]
  \draw[flecha] (0,0) -- (5,0) node[below]{$\ln r$};
  \draw[flecha] (0,0) -- (0,3.2) node[left]{$\Delta h$};
  \draw[dashed] (0,2.6) node[left]{$h_0$} -- (0.8,2.6) -- (0.8,0) node[below]{$r_0$};
  \draw[rojo,thick] (0.8,2.6) -- (3.2,0) node[below,black]{$R$};
  \draw[azulusm,thick] (0.8,2.6) .. controls (1.5,1.6) and (2.8,0.3) .. (3.3,0.05) -- (4.5,0.02);
\end{tikzpicture}
\caption{Estimación de $R$ a partir de mediciones de depresión.}
\end{figure}

\paragraph{Fórmula de Sichardt (1928):} depende de $K$ y, por ende, de la granulometría.
\begin{formula}
$R=3000\,\Delta h_0\sqrt{K}$ \ (m),\qquad $K$ en [m/s],\ $\Delta h_0$ en [m]
\end{formula}
\begin{table}[H]
\centering
\begin{tabular}{lcc}
\toprule
\textbf{Material} & $D$ (mm) & $R$ (m)\\
\midrule
Gravas & mayor a 10 & 1500\\
Gravillas & 2 -- 10 & 500 -- 1500\\
Gravilla-arcilla & 1 -- 2 & 400 -- 500\\
Arena gruesa & 0.5 -- 1 & 200 -- 400\\
Arena media & 0.25 -- 0.5 & 100 -- 200\\
Arena fina & 0.1 -- 0.25 & 50 -- 100\\
Arena muy fina & 0.05 -- 0.1 & 10 -- 50\\
\bottomrule
\end{tabular}
\caption{Valores referenciales del radio de influencia según granulometría.}
\end{table}

\paragraph{Otros métodos empíricos} (ver Louwyck et al., 2022, \emph{The
Radius of Influence Myth}, Water 14(2):149,
\url{https://www.mdpi.com/2073-4441/14/2/149}):
\begin{table}[H]
\centering\small
\begin{tabular}{llll}
\toprule
\textbf{Modelo} & \textbf{Régimen} & \textbf{Borde superior} & \textbf{Radio de influencia $R$}\\
\midrule
Dupuit & Permanente & Nivel freático & Borde exterior (dato)\\
Thiem & Permanente & Impermeable & Borde exterior (dato)\\
de Glee & Permanente & Filtrante & $R=4\sqrt{cKD}$\\
Theis & Transiente & Impermeable & $R=1.5\sqrt{KDt/S}$\\
Hantush--Jacob & Transiente & Filtrante & $R=1.5\sqrt{KDt/S}$ si $t<0.01Sc$;\ $R=4\sqrt{cKD}$ si $t>10Sc$\\
Ernst & Permanente & Drenaje + recarga & $R=\sqrt{Q/(\pi N)}$ o $R=4\sqrt{cKD}$\\
\bottomrule
\end{tabular}
\end{table}
con $KD$: transmisibilidad; $c$: resistencia; $S$: coeficiente de
almacenamiento; $N$: flujo de infiltración; $Q$: caudal de bombeo; $t$: tiempo.
